\documentclass[11pt]{article}
\usepackage[margin=1.1in]{geometry}
\usepackage{amsmath,amssymb,amsthm}
\usepackage{bm}
\usepackage{hyperref}
\usepackage{microtype}
\usepackage{booktabs}

\newtheorem{lemma}{Lemma}
\newtheorem{remark}{Remark}

\newcommand{\nn}{\hat{\bm{n}}}
\newcommand{\mm}{\hat{\bm{m}}}
\newcommand{\qq}{\hat{\bm{q}}}
\newcommand{\ee}[1]{\hat{\bm{e}}_{#1}}
\newcommand{\tL}{\widetilde{L}}
\newcommand{\SD}{\Sigma_D}
\newcommand{\SL}{\Sigma_L}
\newcommand{\sbar}{\bar{\sigma}}
\newcommand{\stilde}{\widetilde{\sigma}}
\newcommand{\invs}{\langle\sigma^{-1}\rangle_{\mathrm{vol}}}
\newcommand{\tvnabla}{\widetilde{\nabla}}
% --- anisotropic extension ---
\newcommand{\snn}[1]{\sigma_{nn}^{(#1)}}
\newcommand{\smn}[1]{\sigma_{mn}^{(#1)}}
\newcommand{\sqn}[1]{\sigma_{qn}^{(#1)}}
\newcommand{\Sc}[1]{S^{(#1)}}
\newcommand{\GTT}{G_{TT}}
\newcommand{\Gnn}{G_{nn}}
\newcommand{\rhom}{\bm{\rho}_m}
\newcommand{\rhoq}{\bm{\rho}_q}
\newcommand{\Done}{\bm{D}_1}

\title{Three-Dimensional Nodal Homogenization for\\
       Arbitrarily Oriented and Multi-Interface Layered Media}
\author{S.\ Moskow}
\date{\today}

\begin{document}
\maketitle

\begin{abstract}
We derive the nodal homogenization tensor $\SD$ for a layered medium
with planar symmetry direction $\nn$ at arbitrary orientation, and describe
its implementation for Cartesian finite-difference grids.
Both the isotropic and fully anisotropic cases are treated within the same
energy-matching framework of Moskow et al.\ (1999), extended to three spatial
dimensions.  In both cases the result takes the form
$\SD = \tL^{-\top} G\,\tL^{-1}$,
where $\tL$ encodes the discrete gradient of the local solution space and $G$
is an energy matrix built from volume averages of the medium.
For an isotropic profile the ingredients are
$\tL = [\mm\mid\qq\mid D\nn]$, with $D$ a diagonal matrix of per-axis line
averages of $\sigma^{-1}$, and $G = \operatorname{diag}(\sbar,\sbar,\invs)$.
For a tensor-valued profile, the columns of $\tL$ acquire off-diagonal
corrections from the off-normal conductivity components, and the upper-left
block of $G$ is replaced by the volume average of the pointwise Schur
complement while $G_{33}$ becomes the volume average of $1/\sigma_{nn}$.

Two practical computational paths are described.  When the interface geometry
is known analytically, $\nn$ is supplied exactly and averages are read from
geometric fractions.  When only a conductivity callable is available, $\nn$
is estimated by fitting a plane to interface-crossing voxel midpoints on a
fine sub-grid (geometric SVD); a planarity ratio detects unreliable estimates
and triggers a diagonal fallback.  For tensor-valued callables the SVD uses
$\tfrac{1}{3}\mathrm{tr}\,\sigma$ as a scalar proxy.

Dual cells straddling two interfaces are handled by sequential nodal homogenization: the two outer materials are homogenized first with the
outer boundary normal, then the result is homogenized against the sandwiched
inner material (identified by its median conductivity) with the inner boundary
normal.  A fallback to single-interface treatment is applied when the
intermediate tensor has eigenvalues outside the physical conductivity range.
\end{abstract}

\section{Setup}

Let $H = [x_1^-,x_1^+]\times[x_2^-,x_2^+]\times[x_3^-,x_3^+]$ be a rectangular
dual cell (node-centered box) on a Cartesian grid with grid axes
$\ee{1},\ee{2},\ee{3}$.  We consider an isotropic medium whose conductivity
varies only along a fixed direction $\nn$:
\begin{equation}\label{eq:layered}
    \sigma(\bm{x}) = \sigma(\nn\cdot\bm{x}),
\end{equation}
where $\nn$ is a unit vector at arbitrary orientation and $\sigma:\mathbb{R}\to\mathbb{R}_{>0}$
is any measurable, bounded, and bounded-away-from-zero function.  This
encompasses any number of planar layers (not just one interface), as well as
continuously varying profiles.

Let $\mm$ and $\qq$ be unit vectors orthogonal to $\nn$ and to each other, so
that $\{\mm,\qq,\nn\}$ is an orthonormal basis.

\section{Local Solution Space}

Following Moskow et al.\ (1999), we approximate the electric potential locally
by functions whose gradient satisfies the interface transmission conditions
exactly at every level set of $\nn\cdot\bm{x}$: the tangential electric field
and the normal current density are both constant in the $\nn$ direction.
This motivates the local solution space
\[
    L(H) = \operatorname{span}\bigl\{1,\;\mm\cdot\bm{x},\;\qq\cdot\bm{x},\;\varphi_n(\bm{x})\bigr\},
\]
where
\[
    \varphi_n(\bm{x}) = \int_0^{\nn\cdot\bm{x}} \frac{dt}{\sigma(t)}
\]
is the ``normal potential,'' which carries all the variability of $\sigma$
across layers.  Note that $\varphi_n$ is well-defined for any $\sigma$
satisfying \eqref{eq:layered}, regardless of the number of layers or
smoothness of the profile.

\section{True and Discrete Gradients}

\paragraph{True gradient.}
For $\phi\in L(H)$ with coefficient vector $\bm{a}=(a_1,a_2,a_3)^\top$,
\[
    \phi(\bm{x}) = a_1(\mm\cdot\bm{x}) + a_2(\qq\cdot\bm{x}) + a_3\,\varphi_n(\bm{x}),
\]
the true gradient is
\begin{equation}\label{eq:true_grad}
    \nabla\phi = a_1\,\mm + a_2\,\qq + \frac{a_3}{\sigma(\nn\cdot\bm{x})}\,\nn.
\end{equation}
The three coefficients have clear physical meaning: $a_1$ and $a_2$ are the
(constant) tangential electric field components, and $a_3$ is the (constant)
normal current density $J_n = \sigma E_n$.

\paragraph{Discrete gradient.}
The finite-difference approximation to the $k$-th component of $\nabla\phi$
along grid axis $\ee{k}$ is the average of $\partial\phi/\partial x_k$ across
the corresponding edge.  Since $\phi = a_1(\mm\cdot\bm{x}) + a_2(\qq\cdot\bm{x})
+ a_3\varphi_n$ and the tangential terms contribute constants $a_1 \hat{m}_k$
and $a_2 \hat{q}_k$, only the normal potential needs care:
\[
    \frac{1}{\Delta x_k}\int_{x_k^-}^{x_k^+}
    \frac{\partial\varphi_n}{\partial x_k}\,dx_k
    = \hat{n}_k \cdot \frac{1}{\Delta x_k}\int_{x_k^-}^{x_k^+}
    \frac{dx_k}{\sigma(\nn\cdot\bm{x})}\Bigr|_{\text{other coords fixed}}
    = \hat{n}_k\, d_k,
\]
where
\begin{equation}\label{eq:dk}
    d_k = \left\langle\sigma^{-1}\right\rangle_k
    \;:=\;
    \frac{1}{\Delta x_k}\int_{x_k^-}^{x_k^+}
    \frac{dx_k}{\sigma(\nn\cdot\bm{x})}\Bigr|_{\text{other coords at node}}
\end{equation}
is the line average of $\sigma^{-1}$ along axis $\ee{k}$ through the node.
This integral is taken at fixed coordinates perpendicular to $\ee{k}$ and
is well-defined for any profile $\sigma(\nn\cdot\bm{x})$.  In the special
case of $M$ planar layers with conductivities $\sigma_1,\ldots,\sigma_M$ and
fractional lengths $f_k^{(j)}$ along edge $k$, it reduces to
$d_k = \sum_j f_k^{(j)}/\sigma_j$.

Assembling all three components, the discrete gradient is
\begin{equation}\label{eq:disc_grad}
    \tvnabla\phi = a_1\,\mm + a_2\,\qq + a_3\,D\nn,
    \qquad D = \operatorname{diag}(d_1,d_2,d_3),
\end{equation}
or in matrix form
\[
    \tvnabla\phi = \tL\,\bm{a},
    \qquad \tL := \bigl[\mm\;\big|\;\qq\;\big|\;D\nn\bigr] \in \mathbb{R}^{3\times 3},
\]
where the columns of $\tL$ are $\mm$, $\qq$, and $D\nn$.

\section{Energy Matching Condition}\label{sec:energy_matching}

We seek an effective tensor $\SD$ such that, for every $\phi,\psi\in L(H)$,
\begin{equation}\label{eq:energy_match}
    |H|\;\tvnabla\phi\cdot\SD\,\tvnabla\psi
    \;=\;
    \int_H \nabla\phi\cdot\sigma\,\nabla\psi\;dV.
\end{equation}

\paragraph{Right-hand side.}
Using \eqref{eq:true_grad} and the orthonormality of $\{\mm,\qq,\nn\}$,
\begin{align*}
    \nabla\phi\cdot\sigma\nabla\psi
    &= \sigma\Bigl(a_1 b_1 + a_2 b_2\Bigr)
     + \frac{a_3 b_3}{\sigma},
\end{align*}
so
\[
    \int_H \nabla\phi\cdot\sigma\nabla\psi\;dV
    = |H|\bigl(\sbar\,(a_1 b_1+a_2 b_2) + \invs\, a_3 b_3\bigr)
    = |H|\;\bm{a}^\top G\,\bm{b},
\]
where
\begin{equation}\label{eq:G}
    G = \operatorname{diag}\!\bigl(\sbar,\;\sbar,\;\invs\bigr),
\end{equation}
with
\[
    \sbar = \langle\sigma\rangle_{\mathrm{vol}} := \frac{1}{|H|}\int_H \sigma(\nn\cdot\bm{x})\,dV,
    \qquad
    \invs = \langle\sigma^{-1}\rangle_{\mathrm{vol}} := \frac{1}{|H|}\int_H \frac{dV}{\sigma(\nn\cdot\bm{x})},
\]
the volume averages of $\sigma$ and $\sigma^{-1}$ over $H$, respectively.
Both are well-defined for any measurable bounded profile \eqref{eq:layered}.
Note that $\invs = 1/\stilde$, where $\stilde = (\invs)^{-1}$ is the harmonic
mean of $\sigma$.  The key observation is that $G_{33} = \invs$ is the
\emph{arithmetic} mean of the resistivity---not the harmonic mean of
$\sigma$---a distinction that is critical when $D\nn \not\propto \nn$.

\paragraph{Left-hand side.}
With $\tvnabla\phi = \tL\bm{a}$ and $\tvnabla\psi = \tL\bm{b}$,
\[
    |H|\;\tvnabla\phi\cdot\SD\,\tvnabla\psi
    = |H|\;\bm{a}^\top\tL^\top\SD\,\tL\,\bm{b}.
\]

\paragraph{Condition.}
Equating both sides for all $\bm{a},\bm{b}$ gives
\[
    \tL^\top\SD\,\tL = G,
\]
and solving for $\SD$:
\begin{equation}\label{eq:SigmaD}
    \boxed{
    \SD = \tL^{-\top}\,G\,\tL^{-1},
    \qquad
    \tL = [\mm\mid\qq\mid D\nn],
    \qquad
    G = \operatorname{diag}\!\bigl(\sbar,\;\sbar,\;\invs\bigr).
    }
\end{equation}
Here $\tL^{-\top} := (\tL^\top)^{-1} = (\tL^{-1})^\top$.

\section{Verification and Special Cases}

\begin{lemma}[Axis-aligned interface]\label{lem:axis}
For $\nn = \ee{3}$, the formula \eqref{eq:SigmaD} reduces to the standard
arithmetic/harmonic tensor
\[
    \SD = \operatorname{diag}(\sbar,\;\sbar,\;\stilde).
\]
\end{lemma}
\begin{proof}
Choose $\mm=\ee{1}$, $\qq=\ee{2}$.  Then $D\nn = d_3\ee{3}$ and
$\tL = \operatorname{diag}(1,1,d_3)$.  Since $\nn = \ee{3}$, lines along
$\ee{1}$ and $\ee{2}$ are parallel to all level sets of $\sigma$, so
$d_1 = d_2 = \langle\sigma^{-1}\rangle_1 = \langle\sigma^{-1}\rangle_2$
each equal the pointwise value $\sigma(\bm{x}_{\mathrm{node}})^{-1}$ (no
variation along those edges).  The $\ee{3}$ edge varies through the full
profile, giving $d_3 = \langle\sigma^{-1}\rangle_z = \invs = 1/\stilde$
(the line average along $z$ equals the volume average when the profile
depends only on $z$).  Therefore
\[
    \tL^{-1} = \operatorname{diag}(1,1,1/d_3) = \operatorname{diag}(1,1,\stilde),
\]
and
\[
    \SD = \operatorname{diag}(1,1,\stilde)\cdot
          \operatorname{diag}(\sbar,\sbar,1/\stilde)\cdot
          \operatorname{diag}(1,1,\stilde)
        = \operatorname{diag}(\sbar,\;\sbar,\;\stilde).\qedhere
\]
\end{proof}

\begin{lemma}[Symmetric straddling]\label{lem:sym}
If $D\nn$ is parallel to $\nn$ (i.e.\ $D = d\,I$ for some scalar $d$), then
$\SD = \SL := \sbar(I - \nn\nn^\top) + \stilde\,\nn\nn^\top$.
\end{lemma}
\begin{proof}
When $D\nn = d\,\nn$ the columns of $\tL$ are $\mm$, $\qq$, $d\nn$---an
orthogonal set with squared norms $1,1,d^2$.  Hence
$\tL^{-1}$ has rows $\mm^\top$, $\qq^\top$, $\nn^\top/d$, giving
$\tL^{-\top} = [\mm\mid\qq\mid\nn/d]$.  Then
\begin{align*}
    \SD &= [\mm\mid\qq\mid\nn/d]\,\operatorname{diag}(\sbar,\sbar,1/\stilde)
           \begin{pmatrix}\mm^\top\\\qq^\top\\\nn^\top/d\end{pmatrix}
    \\
    &= \sbar\,\mm\mm^\top + \sbar\,\qq\qq^\top + \frac{1}{\stilde d^2}\,\nn\nn^\top.
\end{align*}
Since $d = \invs^{1/1} = 1/\stilde$ (the line average equals the volume average
when $D\propto I$), the last coefficient equals $\stilde$, and
$\mm\mm^\top + \qq\qq^\top = I - \nn\nn^\top$, giving $\SD = \SL$.
\end{proof}

\begin{remark}
Lemma~\ref{lem:sym} shows that the nodal correction is zero for any node at
which $D\nn\propto\nn$---equivalently, when $d_x\hat{n}_x : d_z\hat{n}_z$ is
constant along the direction of $\nn$.  For a tilted interface
$\nn=[1,0,1]/\sqrt{2}$, this occurs when $d_x = d_z$, i.e.\ when the edge
fractions in the $x$- and $z$-directions are equal.  The correction is
non-trivial only for nodes where the interface cuts those edges asymmetrically.
\end{remark}

\section{Relation to the Two-Dimensional Formula}

Moskow et al.\ (1999) derive, for the two-dimensional isotropic case with
normal $\nn$ and single tangential $\mm$, the tensor
\[
    \SD^{(2)} = K\,\SL^{(2)}\,K^\top,\qquad K = L^{-1},
\]
where $L$ is a $2\times2$ matrix (their eq.~(A.21)) and
$\SL^{(2)} = \operatorname{diag}(\sbar,\stilde)$ in the $(\mm,\nn)$ basis.

In 3D the same energy-matching argument gives an identical structure.  The
difference is that in 3D there are \emph{two} tangential directions, each
contributing an $\sbar$ term to the diagonal of $G$.  Writing $\SL$ for the
standard $3\times3$ arithmetic/harmonic tensor and noting that
$G = \tL^\top\SL\,\tL$ holds only when $\tL$ is orthogonal (i.e.\ $D\nn \propto
\nn$), one sees that the 3D formula \eqref{eq:SigmaD} is not generally equivalent
to $K\SL K^\top$ with $K=\tL^{-1}$---it is the direct energy-matching result,
which happens to coincide with $K\SL K^\top$ only in the symmetric case.

% ============================================================
\section{Anisotropic Extension}\label{sec:anisotropic}
% ============================================================

We now extend the derivation to the case where the conductivity varies along
$\nn$ as a \emph{tensor-valued} profile: $\sigma = \sigma(\nn\cdot\bm{x})$,
where for each $t\in\mathbb{R}$ the value $\sigma(t)$ is a symmetric
positive-definite matrix.  This includes any number of planar layers as well
as continuously varying anisotropy.  We write
\[
  \sigma_{mn}(t) := \mm^\top\!\sigma(t)\nn,
  \quad
  \sigma_{qn}(t) := \qq^\top\!\sigma(t)\nn,
  \quad
  \sigma_{nn}(t) := \nn^\top\!\sigma(t)\nn > 0,
\]
and use $t = \nn\cdot\bm{x}$ as the layering coordinate.

This section follows the same structure as Appendix~A of Moskow et al.\
(1999), extended from 2D to 3D by the second tangential direction $\qq$.

\subsection*{A.1\quad Interface constants}

For any layered anisotropic medium of the form $\sigma(\nn\cdot\bm{x})$,
the interface transmission conditions at every level set $\{\nn\cdot\bm{x}=t\}$
are continuity of the tangential electric field and of the normal current density.
These motivate three constants that are uniform across all layers:
\begin{align}
  c_1 &:= \mm\cdot\nabla u \quad \text{(tangential field, $\mm$ direction)},
  \label{eq:A1}\\
  c_2 &:= \qq\cdot\nabla u \quad \text{(tangential field, $\qq$ direction)},
  \label{eq:A2}\\
  c_3 &:= \nn\cdot(\sigma\nabla u) \quad \text{(normal current density $J_n$)}.
  \label{eq:A3}
\end{align}

\subsection*{A.2\quad Electric field at each level}

At each level $t$, the condition $c_3 = \sigma_{mn}(t)c_1 + \sigma_{qn}(t)c_2
+ \sigma_{nn}(t)E_n$ determines the normal field:
\begin{equation}\label{eq:A4}
  E_n(t) \;=\; \frac{c_3 - \sigma_{mn}(t)\,c_1 - \sigma_{qn}(t)\,c_2}{\sigma_{nn}(t)}.
\end{equation}
The full gradient at level $t$ is therefore
\begin{equation}\label{eq:A5}
  \nabla u\big|_{t}
  \;=\; c_1\,\mm \;+\; c_2\,\qq \;+\; E_n(t)\,\nn,
\end{equation}
and the triple $(c_1,c_2,c_3)$ completely characterizes the local potential.

\subsection*{A.3\quad Discrete gradient}

The finite-difference approximation to $(\nabla u)_k$ along grid axis $k$
averages \eqref{eq:A5} over the corresponding edge at fixed transverse
coordinates.  Because $c_1,c_2$ are constant, only $E_n(t)$ varies along
the edge.  Substituting \eqref{eq:A4} and collecting terms in $c_1,c_2,c_3$:
\begin{equation}\label{eq:A7}
  (\tvnabla u)_k
  \;=\; c_1\!\left(\hat{m}_k - \hat{n}_k\,[D_2]_{kk}\right)
      + c_2\!\left(\hat{q}_k - \hat{n}_k\,[D_3]_{kk}\right)
      + c_3\,\hat{n}_k\,[D_1]_{kk},
\end{equation}
where the per-axis line averages are
\begin{align}
  [D_1]_{kk} &\;:=\;
    \frac{1}{\Delta x_k}\int_{x_k^-}^{x_k^+}
    \frac{dx_k}{\sigma_{nn}(\nn\cdot\bm{x})}\bigg|_{x_j = x_j^{\mathrm{node}},\,j\neq k},
    \label{eq:A8}\\[4pt]
  [D_2]_{kk} &\;:=\;
    \frac{1}{\Delta x_k}\int_{x_k^-}^{x_k^+}
    \frac{\sigma_{mn}(\nn\cdot\bm{x})}{\sigma_{nn}(\nn\cdot\bm{x})}\,dx_k\bigg|_{x_j = x_j^{\mathrm{node}},\,j\neq k},
    \label{eq:A9}\\[4pt]
  [D_3]_{kk} &\;:=\;
    \frac{1}{\Delta x_k}\int_{x_k^-}^{x_k^+}
    \frac{\sigma_{qn}(\nn\cdot\bm{x})}{\sigma_{nn}(\nn\cdot\bm{x})}\,dx_k\bigg|_{x_j = x_j^{\mathrm{node}},\,j\neq k}.
    \label{eq:A10}
\end{align}
Each is a 1D integral along grid axis $k$ through the node, with the two
transverse coordinates held fixed at their node values.  Because
$\nn\cdot\bm{x} = \hat{n}_k x_k + \sum_{j\neq k}\hat{n}_j x_j^{\mathrm{node}}$,
the integrand is a function of $x_k$ alone, so the integrals are
well-defined for any measurable profile.  This is the exact anisotropic
analogue of $d_k$ in \eqref{eq:dk}.  Thus $D_1$ is the per-axis line
average of $1/\sigma_{nn}$; $D_2$ and $D_3$ are the line averages of the
ratios $\sigma_{mn}/\sigma_{nn}$ and $\sigma_{qn}/\sigma_{nn}$.
Assembling all three Cartesian components:
\begin{equation}\label{eq:A11}
  \tvnabla u \;=\; \tL\,\bm{c},
  \qquad
  \tL \;:=\; \bigl[\,\mm - D_2\nn \;\big|\; \qq - D_3\nn \;\big|\; D_1\nn\,\bigr],
  \qquad
  \bm{c} = (c_1, c_2, c_3)^\top,
\end{equation}
where $D_j\nn$ denotes the vector with $k$-th entry $[D_j]_{kk}\hat{n}_k$.

\begin{remark}[2D reduction]
In 2D with a single tangential direction $\mm$, there is no $c_2$ or $D_3$,
and $D_2$ reduces to the $2\times 2$ diagonal matrix of line averages of
$\sigma_{mn}/\sigma_{nn}$.  Then $\tL = [\mm - D_2\nn \mid D_1\nn]$, matching
eq.~(A.8) of Moskow et al.\ (1999).
\end{remark}

\subsection*{A.4\quad Pointwise Schur complement and energy decomposition}

At each level $t$, define the \emph{pointwise Schur complement} of $\sigma(t)$
with respect to $\nn$:
\begin{equation}\label{eq:A12}
  S(t) \;:=\; \sigma(t)
  \;-\; \frac{(\sigma(t)\nn)\,(\nn^\top\!\sigma(t))}{\sigma_{nn}(t)}.
\end{equation}
$S(t)$ is symmetric positive-semidefinite with null vector $\nn$ for every $t$.
Substituting \eqref{eq:A4}--\eqref{eq:A5} into the pointwise energy density
and using the symmetry of $\sigma(t)$, the normal-tangential cross-terms
cancel and one obtains the clean factorization
\begin{equation}\label{eq:A13}
  \nabla u\cdot\sigma(t)\nabla u
  \;=\; \bm{\xi}^\top\!\bigl(\mm\;\qq\bigr)^\top\!S(t)\bigl(\mm\;\qq\bigr)\bm{\xi}
  \;+\; \frac{c_3^2}{\sigma_{nn}(t)},
  \qquad \bm{\xi} = (c_1,c_2)^\top.
\end{equation}
The tangential and normal contributions are \emph{decoupled} pointwise.
Integrating over $H$:
\begin{equation}\label{eq:A15}
  \frac{1}{|H|}\int_H \nabla u\cdot\sigma\nabla u\;dV
  \;=\; \bm{c}^\top G\,\bm{c},
\end{equation}
where the $3\times 3$ energy matrix in the $(c_1,c_2,c_3)$ basis is
\begin{equation}\label{eq:A16}
  G \;=\;
  \begin{pmatrix}
    \mm^\top\!\GTT\mm & \mm^\top\!\GTT\qq & 0 \\[2pt]
    \qq^\top\!\GTT\mm & \qq^\top\!\GTT\qq & 0 \\[2pt]
    0 & 0 & \Gnn
  \end{pmatrix},
\end{equation}
with volume averages
\begin{equation}\label{eq:GTT_Gnn}
  \GTT \;:=\; \frac{1}{|H|}\int_H S(\nn\cdot\bm{x})\,dV,
  \qquad
  \Gnn \;:=\; \frac{1}{|H|}\int_H \frac{dV}{\sigma_{nn}(\nn\cdot\bm{x})}.
\end{equation}
That is, $\GTT$ is the \emph{arithmetic} volume average of the pointwise
Schur complement, and $\Gnn$ is the \emph{arithmetic} volume average of the
normal resistivity $1/\sigma_{nn}$---the anisotropic analogue of $\invs$ in the
isotropic case.
The zeros in \eqref{eq:A16} follow from $S(t)\nn = 0$ for all $t$, which gives
$\GTT\nn = 0$.

\subsection*{A.5\quad Effective tensor}

The energy-matching condition
$\tL^\top\SD\,\tL = G$ (Section~\ref{sec:energy_matching}) gives
\begin{equation}\label{eq:A17}
  \boxed{
  \SD \;=\; \tL^{-\top}\,G\,\tL^{-1},
  }
\end{equation}
with $\tL$ and $G$ defined by \eqref{eq:A11} and \eqref{eq:A16}--\eqref{eq:GTT_Gnn}.
Equivalently, writing $K = \tL^{-1}$,
\begin{equation}\label{eq:A18}
  \SD \;=\; K^\top\,G\,K,
\end{equation}
the exact 3D analogue of eq.~(A.14) in Moskow et al.\ (1999).  All
ingredients ($D_1,D_2,D_3$ and the averages $\GTT,\Gnn$) are well-defined
for any measurable tensor-valued profile $\sigma(\nn\cdot\bm{x})$.

\begin{remark}[Piecewise-constant reduction]
For a two-layer medium with $\sigma^{(1)},\sigma^{(2)}$ and volume fractions
$f$, $1-f$ (and line fractions $f_k$, $1-f_k$), the integrals reduce to
\begin{align*}
  [D_1]_{kk} &= \frac{f_k}{\sigma_{nn}^{(1)}} + \frac{1-f_k}{\sigma_{nn}^{(2)}},
  &\qquad
  \GTT &= f\,S^{(1)} + (1-f)\,S^{(2)},\\
  [D_2]_{kk} &= \frac{f_k\,\sigma_{mn}^{(1)}}{\sigma_{nn}^{(1)}} + \frac{(1-f_k)\,\sigma_{mn}^{(2)}}{\sigma_{nn}^{(2)}},
  &\qquad
  \Gnn &= \frac{f}{\sigma_{nn}^{(1)}} + \frac{1-f}{\sigma_{nn}^{(2)}},
\end{align*}
and similarly for $D_3$.
\end{remark}

\begin{remark}[Isotropic reduction]
When $\sigma(t)=\sigma(t)\,I$ is scalar, $\sigma_{mn}=\sigma_{qn}=0$ so
$D_2=D_3=0$.  Then $\tL=[\mm\mid\qq\mid D_1\nn]$ with $[D_1]_{kk}=d_k$
the isotropic line average, recovering \eqref{eq:disc_grad}.  The Schur
complement is $S(t)=\sigma(t)(I-\nn\nn^\top)$, giving $\GTT=\sbar(I-\nn\nn^\top)$
and $\Gnn=\invs$, so $G=\operatorname{diag}(\sbar,\sbar,\invs)$, recovering \eqref{eq:G}.
\end{remark}

\begin{remark}[TI medium at oblique interface]
For a transversely isotropic (TI) medium $\sigma=\sigma_T(I-\hat{\bm{t}}\hat{\bm{t}}^\top)
+\sigma_N\hat{\bm{t}}\hat{\bm{t}}^\top$ with symmetry axis $\hat{\bm{t}}$,
and with interface normal $\nn$ not aligned with $\hat{\bm{t}}$, the
ratios $\sigma_{mn}/\sigma_{nn}$ and $\sigma_{qn}/\sigma_{nn}$ are nonzero,
making $D_2$ and $D_3$ nontrivial.  Omitting them leads to a $k$-dependent
crossing position in borehole logging simulations.
\end{remark}

% ============================================================
\section{Summary}

The three-dimensional nodal homogenization tensor for a planar interface with
unit normal $\nn$ is given by
\[
    \SD = \tL^{-\top}\,G\,\tL^{-1},
\]
with the following ingredients:

\medskip
\begin{center}
\renewcommand{\arraystretch}{1.6}
\begin{tabular}{lll}
\toprule
Quantity & Definition & Meaning \\
\midrule
$D$ & $\operatorname{diag}(d_1,d_2,d_3)$,\quad
      $d_k = \dfrac{1}{\Delta x_k}\displaystyle\int_{x_k^-}^{x_k^+}
      \dfrac{dx_k}{\sigma(\nn\cdot\bm{x})}$
    & per-axis line avg of $\sigma^{-1}$ \\[4pt]
$\tL$ & $[\mm\mid\qq\mid D\nn]$, \quad $\mm,\qq\perp\nn$ orthonormal & discrete-gradient map \\
$G$ & $\operatorname{diag}(\sbar,\sbar,\invs)$ & energy matrix \\
$\sbar$ & $\dfrac{1}{|H|}\displaystyle\int_H \sigma(\nn\cdot\bm{x})\,dV$ & vol.\ avg.\ of $\sigma$ \\[4pt]
$\invs$ & $\dfrac{1}{|H|}\displaystyle\int_H \dfrac{dV}{\sigma(\nn\cdot\bm{x})}$ & vol.\ avg.\ of $\sigma^{-1}$ \\
\bottomrule
\end{tabular}
\end{center}

\medskip\noindent
All quantities are well-defined for any measurable profile
$\sigma = \sigma(\nn\cdot\bm{x})$, with any number of layers or continuous
variation.  The key point is that $G_{33} = \invs$ is the
\emph{arithmetic mean of the resistivity} over the dual-cell volume, not
the harmonic mean of $\sigma$.  This value arises from the integral
$\int_H \sigma\,|E_n|^2\,dV = \int_H \sigma^{-1}\,|J_n|^2\,dV$
when $(a_1,a_2,a_3)=(0,0,1)$.

\subsection*{Anisotropic case}

For a tensor-valued layered profile $\sigma = \sigma(\nn\cdot\bm{x})$, the
same formula $\SD = \tL^{-\top}G\,\tL^{-1}$ holds with the following
ingredients (see Section~\ref{sec:anisotropic}):

\medskip
\begin{center}
\renewcommand{\arraystretch}{1.8}
\begin{tabular}{lll}
\toprule
Quantity & Definition & Meaning \\
\midrule
$D_1$ & $\operatorname{diag}_k\!\left(\dfrac{1}{\Delta x_k}
        \displaystyle\int_{x_k^-}^{x_k^+}
        \dfrac{dx_k}{\sigma_{nn}(\nn\cdot\bm{x})}\right)$
      & line avg of $\sigma_{nn}^{-1}$ \\[6pt]
$D_2$ & $\operatorname{diag}_k\!\left(\dfrac{1}{\Delta x_k}
        \displaystyle\int_{x_k^-}^{x_k^+}
        \dfrac{\sigma_{mn}}{\sigma_{nn}}\,dx_k\right)$
      & line avg of $\sigma_{mn}/\sigma_{nn}$ \\[6pt]
$D_3$ & $\operatorname{diag}_k\!\left(\dfrac{1}{\Delta x_k}
        \displaystyle\int_{x_k^-}^{x_k^+}
        \dfrac{\sigma_{qn}}{\sigma_{nn}}\,dx_k\right)$
      & line avg of $\sigma_{qn}/\sigma_{nn}$ \\[6pt]
$\tL$ & $[\mm - D_2\nn \mid \qq - D_3\nn \mid D_1\nn]$
      & discrete-gradient map \\[2pt]
$S(t)$ & $\sigma(t) - \dfrac{(\sigma(t)\nn)(\nn^\top\!\sigma(t))}{\sigma_{nn}(t)}$
       & pointwise Schur complement \\[6pt]
$\GTT$ & $\dfrac{1}{|H|}\displaystyle\int_H S(\nn\cdot\bm{x})\,dV$
       & vol.\ avg.\ of Schur complement \\[6pt]
$\Gnn$ & $\dfrac{1}{|H|}\displaystyle\int_H \dfrac{dV}{\sigma_{nn}(\nn\cdot\bm{x})}$
       & vol.\ avg.\ of $\sigma_{nn}^{-1}$ \\[6pt]
$G$ & $\begin{pmatrix} \mm^\top\!\GTT\mm & \mm^\top\!\GTT\qq & 0 \\
                       \qq^\top\!\GTT\mm & \qq^\top\!\GTT\qq & 0 \\
                       0 & 0 & \Gnn \end{pmatrix}$
    & energy matrix \\
\bottomrule
\end{tabular}
\end{center}

\medskip\noindent
The isotropic case is recovered when $\sigma(t) = \sigma(t)I$: then
$D_2 = D_3 = 0$, $\tL$ reduces to $[\mm\mid\qq\mid D\nn]$ with
$D = D_1$, and $G = \operatorname{diag}(\sbar,\sbar,\invs)$.

% ============================================================
\section{Practical Application: Two Computational Paths}
% ============================================================

The derivation above assumes $\sigma$ is exactly layered with known normal
$\nn$ inside the dual cell $H$.  In practice two situations arise, each handled
by a distinct computational path.

% ------------------------------------------------------------
\subsection{Case 1 — Interface geometry analytically known}\label{sec:exact_n}
% ------------------------------------------------------------

When the interface is described by an explicit geometric object (e.g.\ a
planar boundary, a cylinder, a level-set surface), the unit normal $\nn$ at
each dual-cell node is available analytically.  The procedure is then direct:

\begin{enumerate}
\item \textbf{Evaluate material parameters.}  Evaluate $\sigma$ (scalar or
tensor) on a fine isotropic sub-grid of spacing $h_\mathrm{sub}$ within $H$,
using the exact interface geometry to assign the correct material to each
sub-voxel.
\item \textbf{Volume fractions.}  From the sub-grid masks, compute the
volume fraction $f$ of each material in $H$ and the per-axis line fractions
$f_k$ along the central grid-axis lines.
\item \textbf{Line averages.}  The per-axis line averages $d_k$
(eq.\ \eqref{eq:dk}) are read directly from $f_k$:
\[
  d_k \;=\; \frac{f_k}{\sigma_1} + \frac{1-f_k}{\sigma_2}.
\]
For the anisotropic case replace $\sigma_i$ by the appropriate components
$\sigma_{nn}^{(i)}$, $\sigma_{mn}^{(i)}/\sigma_{nn}^{(i)}$, etc.
\item \textbf{Volume averages.}  Compute $\sbar = f\sigma_1 + (1-f)\sigma_2$
and $\invs = f/\sigma_1 + (1-f)/\sigma_2$ from the volume fractions.
\item \textbf{Assemble $\SD$.}  Apply \eqref{eq:SigmaD} with the known $\nn$
and the computed ingredients.
\end{enumerate}

No interface-direction estimation is needed; the analytic $\nn$ is used
exactly.  This path is exact for planar interfaces and $O(h^2)$-accurate
for curved ones (where the local planarity error over a cell of width $h$
is $O(h^2)$).

% ------------------------------------------------------------
\subsection{Case 2 — Only $\sigma$ available as a lookup function}
\label{sec:sigma_func}
% ------------------------------------------------------------

When $\sigma(\bm{x})$ is given only as a black-box callable (e.g.\ a
nearest-neighbour interpolation of a fine volumetric data set), no explicit
interface geometry is available.  The direction $\nn$ must be \emph{estimated}
from the local variation of $\sigma$ within $H$, and the line and volume
averages computed entirely by quadrature.

\subsubsection*{Step 1 — Sub-grid evaluation}

Evaluate $\sigma$ on a fine isotropic sub-grid of spacing $h_\mathrm{svd}$
within $H$, producing a 3-D array $B_{ijk} = \sigma(\bm{x}_{ijk})$.  Equal
spacing in all three directions is important: it eliminates sampling-induced
bias in the SVD step below.  The same block is reused for normal estimation,
volume averages, and line averages.

\subsubsection*{Step 2 — Normal estimation by SVD}

When $\sigma$ is tensor-valued, the scalar proxy $b_{ijk} = \tfrac{1}{3}\mathrm{tr}\,B_{ijk}$
is used in place of $B_{ijk}$ for all steps below; the actual tensor values are
retained in a parallel block for use in the homogenization integrals (Steps 3--5).
This ensures that the interface-detection logic is insensitive to the anisotropy
\emph{direction} and responds only to changes in conductivity \emph{magnitude}.

The interface normal is estimated geometrically from the positions of
interface-crossing voxel pairs, not from finite differences of $\sigma$.
(A gradient-based approach is biased when the grid spacing is anisotropic:
a 7:1 $z/x$ ratio inflates $\partial\sigma/\partial z$ by $7\times$ relative
to $\partial\sigma/\partial x$, producing a spurious normal near $\hat{z}$.)

\begin{enumerate}
\item \textbf{Binarise.}  Set $\sigma_\mathrm{mid} = \tfrac{1}{2}(\min b + \max b)$;
      label each sub-voxel as layer~1 ($b < \sigma_\mathrm{mid}$) or
      layer~2 ($b \geq \sigma_\mathrm{mid}$).
\item \textbf{Locate crossings.}  For each of the six face-adjacency directions
      $(\pm x, \pm y, \pm z)$, flag voxel pairs whose two members belong to
      different layers.  Record the physical midpoint of each such pair.
      (Using the midpoint rather than a single voxel centre places each sample
      directly on the interface and avoids a ``flooding'' bias near corners.)
\item \textbf{Fit a plane.}  Stack the midpoint coordinates into an $N\times 3$
      matrix $P$, remove its centroid, and compute the SVD
      $P - \bar P = U\,S\,V^\top$.  The \emph{last} right singular vector
      $V_{:,3}$ (minimum-variance direction of the midpoints) is the estimated
      plane normal $\nn$.
\item \textbf{Sign convention.}  Orient $\nn$ so that the half-space
      $\{\bm{x}\cdot\nn > 0\}$ has higher mean conductivity.
\end{enumerate}

The confidence of the estimate is measured by the \emph{planarity ratio}
\[
  r = \frac{s_3}{s_1}
\]
(last to first singular value of $P - \bar P$).  For a perfectly planar
interface $s_3 \approx 0$ so $r \approx 0$; for an isotropic scatter
$r \approx 1$.  If $r \geq \tau$ (default $\tau = 0.7$) the block has no
single dominant interface direction and the cell falls back to the diagonal
(axis-aligned arithmetic/harmonic) formula.  Otherwise $\nn$ is accepted and
the nodal formula \eqref{eq:SigmaD} is applied.

\subsubsection*{Step 3 — Line averages by 1-D quadrature}

Once $\nn$ is known, the per-axis line averages
\[
  d_k = \frac{1}{\Delta x_k}\int_{x_k^-}^{x_k^+} \frac{dx_k}{\sigma(\nn\cdot\bm{x})}
\]
(eq.\ \eqref{eq:dk}) are computed by 1-D numerical quadrature along each
grid axis through the node centre, with $n_\mathrm{line}$ uniformly spaced
points (default 50).  Similarly for the anisotropic integrands
$\sigma_{mn}/\sigma_{nn}$ and $\sigma_{qn}/\sigma_{nn}$.

\subsubsection*{Step 4 — Volume averages by 3-D quadrature}

The volume averages
$\sbar = \langle\sigma\rangle_\mathrm{vol}$ and
$\invs = \langle\sigma^{-1}\rangle_\mathrm{vol}$
are computed from the same isotropic SVD sub-grid (which already provides
a fine, uniform sample of $H$) rather than a separate coarser $n_\mathrm{vol}^3$
grid, making them more accurate for thin-sliver cells.

\subsubsection*{Step 5 — Assemble $\SD$}

With the estimated $\nn$ and numerically computed $D = \operatorname{diag}(d_1,d_2,d_3)$,
$\sbar$, $\invs$, apply \eqref{eq:SigmaD}.  Cells for which the SVD
confidence is low ($r \geq \tau$) or for which $H$ contains no interface
(uniform block) receive the diagonal tensor
$\SD = \operatorname{diag}(\sigma_{xx}^{\mathrm{eff}}, \sigma_{yy}^{\mathrm{eff}}, \sigma_{zz}^{\mathrm{eff}})$
from axis-wise harmonic averaging.

\begin{remark}
The lookup-function path is fully matrix-free: it calls $\sigma(\bm{x})$
on demand at whatever resolution is required for each dual cell, without
storing a global fine-grid array.  For a piecewise-constant closure the
total evaluation count per node is approximately $n_\mathrm{svd}^3 + 3\,n_\mathrm{line}$,
where $n_\mathrm{svd} = \lceil \Delta x / h_\mathrm{svd}\rceil \times \cdots$.
\end{remark}

% ============================================================
\section{Multi-Interface Cells: Sequential Nodal Homogenization}\label{sec:multilayer}
% ============================================================

In practical geometries a dual cell may straddle \emph{two or more} planar
(or approximately planar) interfaces simultaneously.  A common example in
borehole logging is a cell that contains three distinct materials:
borehole fluid (conductivity $\sigma_b$), invasion zone ($\sigma_i$), and
undisturbed formation ($\sigma_f$), separated by the borehole wall
(inner boundary, normal $\nn_\mathrm{in}$) and the invasion front
(outer boundary, normal $\nn_\mathrm{out}$).  A single interface normal
cannot represent both boundaries at once.

\subsection*{Multi-material identification (lookup-function path)}

A cell is classified as multi-material during the same pass that computes the
SVD normal estimate (\S\ref{sec:sigma_func}).  The isotropic SVD sub-grid block
$B_{ijk}$ (already evaluated for normal estimation) is rounded to six significant
figures and its unique conductivity levels are extracted:
\[
  \{\sigma^{(1)} < \sigma^{(2)} < \cdots < \sigma^{(M)}\}
  \;=\; \mathrm{unique}\!\left(\mathrm{round}(B_{ijk},\,6)\right).
\]
If $M \geq 3$ the cell is flagged as multi-material and the procedure below is
applied.  The rounding tolerance (six decimal places) groups physically identical
conductivities that may differ by floating-point noise without merging genuinely
distinct materials whose conductivities differ by more than $10^{-6}$.

\medskip
\noindent\textbf{Determining interface structure via per-pair sub-SVDs.}
For each of the three pairs of adjacent conductivity levels,
$(0,1)$, $(1,2)$, and the \emph{skip pair} $(0,2)$, a masked sub-SVD is
performed on the voxels belonging to that pair, using an adaptive tolerance
$\delta_{ij} = 0.4 \times \min(\text{gap to nearest other level})$ to define
the masks.  Each sub-SVD returns an estimated normal $\hat{n}_{ij}$ and a
planarity ratio $r_{ij}$; a pair is \emph{resolved} if $r_{ij} < \tau$
(same threshold $\tau = 0.7$ as the single-interface case).

The resolved normals determine the cell structure:
\begin{itemize}
\item \textbf{All interfaces parallel}
      ($|\hat{n}_{01}\cdot\hat{n}_{12}| > 0.90$ and
       $|\hat{n}_{01}\cdot\hat{n}_{02}| > 0.90$):
      the three materials are co-planar strata sharing a single normal
      direction.  The nodal formula \eqref{eq:SigmaD} is applied once to all
      three materials simultaneously via a direct multi-region extension
      (\emph{multiregion nodal}), using the average of the resolved normals.
\item \textbf{Adjacent pairs parallel, skip pair distinct}
      ($\hat{n}_{01} \parallel \hat{n}_{12}$, $\hat{n}_{02} \not\parallel \hat{n}_{01}$):
      two interfaces with two distinct normals — sequential nodal homogenization applies
      (see below).  The \emph{primary} normal $\nn_\mathrm{in}$ is the average
      of $\hat{n}_{01}$ and $\hat{n}_{12}$; the \emph{secondary} normal
      $\nn_\mathrm{out}$ is $\hat{n}_{02}$.
\item \textbf{Otherwise}: the structure is ambiguous; the cell falls back to
      the standard single-interface nodal path using the global SVD normal.
\end{itemize}

\subsection*{Sequential nodal homogenization for $M = 3$ materials}

Label the three distinct materials in sorted order with conductivities
$\sigma_1 < \sigma_2 < \sigma_3$ (or mean tensors $\sigma^{(j)}$ in the
anisotropic case) and volume fractions $f_1, f_2, f_3$ in $H$.  Let
$\nn_\mathrm{in}$ and $\nn_\mathrm{out}$ be the inner and outer boundary
normals (innermost boundary first).  The \emph{inner} (sandwiched) material
is always identified as the one with the \emph{middle} conductivity value
$\sigma_2$; the two \emph{outer} materials are $\sigma_1$ and $\sigma_3$.

Both homogenization steps apply the same 2-material nodal formula
\eqref{eq:SigmaD} — there is no switch to a simpler (e.g.\ Backus/laminate)
formula at either level:

\begin{enumerate}
\item \textbf{Outer homogenization.}
Apply the nodal formula \eqref{eq:SigmaD} to the two
\emph{outer} materials $(\sigma_1, f_1)$ and $(\sigma_3, f_3)$ with
normal $\nn_\mathrm{out}$, treating them as if they occupied the whole cell
in proportion $f_1/(f_1+f_3)$ and $f_3/(f_1+f_3)$.  Per-axis line fractions
are computed from the central-axis lines of the sub-grid restricted to the
union of the outer-material voxels:
\[
  \SD^{(13)} \;=\; \tL_{13}^{-\top}\,G_{13}\,\tL_{13}^{-1}.
\]
\item \textbf{Inner homogenization.}
Treat $\SD^{(13)}$ as one effective material and the inner material $\sigma_2$
as the other.  Apply the nodal formula again with normal $\nn_\mathrm{in}$ and
volume fractions $f_2$ (inner) and $f_1+f_3$ (outer):
\[
  \SD \;=\; \tL_\mathrm{in}^{-\top}\,G_\mathrm{in}\,\tL_\mathrm{in}^{-1}.
\]
Per-axis line fractions for step 2 are read from the full cell's central-axis
lines (not restricted to the outer-material sub-block), so they correctly
capture the inner material's extent relative to the whole dual cell.
\end{enumerate}

\subsection*{Interface normals: exact vs.\ estimated}

In the \emph{exact geometry} path (\S\ref{sec:exact_n}), both $\nn_\mathrm{in}$
and $\nn_\mathrm{out}$ are supplied analytically by the interface description
(e.g.\ the cylinder's outward normal and the dip plane's normal at the node).
No sub-SVD analysis is needed.

In the \emph{lookup-function} path (\S\ref{sec:sigma_func}), the normals are
determined by the per-pair sub-SVDs described above.

\subsection*{Fallback}

If the sequential nodal homogenization result has a maximum eigenvalue exceeding three times
the largest material conductivity (a symptom of near-degeneracy when one material
has very low normal conductivity and dominates $G_{nn}$), the algorithm falls
back to the standard single-interface 2-material nodal formula using the
best-separating normal among all candidates.

\begin{remark}[Nesting convention]
The order of the two homogenization steps matters.  The convention ``outer materials
first'' is chosen so that the inner sandwiched material is always
homogenized last, against a background that already reflects the outer
structure.  For DDH03-type geometries (cylindrical borehole inside a dipping
formation), $\nn_\mathrm{in}$ is the cylinder axis normal and
$\nn_\mathrm{out}$ is the formation dip-plane normal.
\end{remark}

\bigskip
\noindent\textbf{Reference.}\\
S.\ Moskow, V.\ Druskin, T.\ Habashy, P.\ Lee, and S.\ Davydycheva,
\emph{A finite difference scheme for elliptic equations with rough coefficients
using a Cartesian grid nonconforming to interfaces},
SIAM J.\ Numer.\ Anal., 36(2):442--464, 1999.

\end{document}
